% ============================================================
% CAPÍTULO 5 — IMPLEMENTAÇÃO DA APLICAÇÃO
% Sistema de Gestão Integrada — Clube Amigos de Polvoreira (CAP)
% ============================================================

\chapter{Implementação da Aplicação}\label{chap:implementacao}

% ============================================================
% SECÇÃO 5.1 — PROCESSO DE IMPLEMENTAÇÃO
% ============================================================

\section{Processo de Implementação}

Esta secção descreve como a equipa se organizou para implementar a aplicação. As secções seguintes apresentam o processo de desenvolvimento de cada camada do sistema.

As tecnologias usadas seguiram as decisões da fase de especificação: \textbf{Monólito Modular} em \textbf{C\# com ASP.NET Core (.NET 7)} \cite{aspnetcore2023}, base de dados \textbf{PostgreSQL 15} \cite{postgresql2024} via \textbf{Entity Framework Core} \cite{efcore2023}, \textbf{API RESTful} e \textit{frontend} em \textbf{Next.js 16} \cite{nextjs2024} (React 19, TypeScript). A equipa usou também \textbf{Docker Compose} \cite{docker2024} para contentorizar a aplicação e \textbf{Git} para controlo de versões.

O primeiro passo foi definir as \textbf{entidades de domínio} em C\# --- \texttt{Utilizador}, \texttt{Atleta}, \texttt{Treinador}, \texttt{Equipa}, \texttt{Treino}, \texttt{AtestadoMedico}, \texttt{Quota}, entre outras. Com a abordagem \textit{Code-First} do Entity Framework Core, as tabelas foram criadas automaticamente a partir dessas classes via \textit{Migrations}. Cada módulo tem o seu próprio \texttt{DbContext} e esquema na base de dados. Em paralelo, foram criados os primeiros \textit{Controllers} para verificar que a comunicação entre a API e a base de dados funcionava.

Com os modelos estabilizados, a equipa desenvolveu a \textbf{camada de negócio} --- os \textit{Controllers} RESTful dos sete módulos. Cada controlador foi testado com o \textbf{Swagger UI}, disponível em \texttt{http://localhost:5000/swagger}, que permitiu validar os \textit{endpoints} antes de avançar para o \textit{frontend}.

Em paralelo, foi desenvolvida a \textbf{camada de apresentação} em Next.js com TypeScript e Tailwind CSS. O \textit{frontend} seguiu uma abordagem incremental: primeiro com dados estáticos para definir o \textit{layout}, depois substituídos por chamadas reais à API via uma função utilitária \texttt{fetchApi}. Os \textit{dashboards} foram organizados por perfil de utilizador.

Com as três camadas funcionais, passou-se à \textbf{integração e povoamento}. Foi criado um componente \texttt{DataSeeder} que povoa a base de dados com dados de demonstração realistas: utilizadores de todos os papéis, equipas, treinos com presenças variadas, atestados com diferentes estados e quotas com pagamentos em vários estados. A aplicação foi depois testada com os diferentes perfis.

Toda a \textit{stack} foi \textbf{contentorizada com Docker}. O \texttt{docker-compose.yml} orquestra três serviços: a base de dados (PostgreSQL~15 Alpine), a API (.NET) e o \textit{frontend} (Next.js). Um único comando --- \texttt{docker compose up -{}-build} --- é suficiente para compilar, aplicar as migrações, povoar os dados e servir a aplicação nos portos 3000 e 5000.

% ============================================================
% SECÇÃO 5.2 — IMPLEMENTAÇÃO DA BASE DE DADOS
% ============================================================

\section{Implementação da Base de Dados}

O projeto usou a abordagem \textit{Code-First} do Entity Framework Core. As tabelas e relações são criadas automaticamente a partir das classes C\#, sem necessidade de escrever SQL manualmente.

\subsection{Criação da Base de Dados e Esquemas}

A base de dados é um único servidor PostgreSQL~15 dividido em sete \textbf{esquemas}, um por módulo, para garantir o isolamento lógico dos dados de cada área. A tabela seguinte resume os esquemas e as tabelas principais de cada um:

\begin{center}
{\footnotesize
\rowcolors{2}{rowblue}{white}
\begin{xltabular}{\textwidth}{|>{\raggedright\arraybackslash}p{2.5cm}|>{\raggedright\arraybackslash}p{3.5cm}|>{\raggedright\arraybackslash}X|}
\hline
\rowcolor{headerblue}\textcolor{white}{\textbf{Esquema}} & \textcolor{white}{\textbf{Módulo}} & \textcolor{white}{\textbf{Tabelas Principais}} \\ \hline
\endfirsthead

\multicolumn{3}{c}{{\tablename\ \thetable{} -- Continuação}} \\
\hline
\rowcolor{headerblue}\textcolor{white}{\textbf{Esquema}} & \textcolor{white}{\textbf{Módulo}} & \textcolor{white}{\textbf{Tabelas Principais}} \\ \hline
\endhead

\hline \multicolumn{3}{|r|}{\textit{Continua na próxima página\ldots}} \\ \hline \endfoot

\hline
\caption{Esquemas da base de dados e respetivas tabelas principais.}
\label{tab:esquemas-bd} \\
\endlastfoot

\ttsplit{users}         & Gestão de Utilizadores     & \ttsplit{Utilizadores} (com TPH para Atleta, Treinador, Encarregado) \\ \hline
\ttsplit{sports}        & Gestão Desportiva          & \ttsplit{Modalidades}, \ttsplit{Escaloes}, \ttsplit{Equipas}, \ttsplit{AtletaEquipas}, \ttsplit{Treinos}, \ttsplit{PresencasTreino}, \ttsplit{Convocatorias}, \ttsplit{Convites}, \ttsplit{OcorrenciasJogo}, \ttsplit{AvaliacoesQualitativas} \\ \hline
\ttsplit{clinical}      & Gestão Clínica             & \ttsplit{Atestados}, \ttsplit{Lesoes} \\ \hline
\ttsplit{finance}       & Gestão Financeira          & \ttsplit{QuotaDefinicoes}, \ttsplit{Quotas}, \ttsplit{Pagamentos}, \ttsplit{Recibos}, \ttsplit{PagamentoQuota} \\ \hline
\ttsplit{notifications} & Serviço de Notificações    & \ttsplit{Notificacoes}, \ttsplit{NotificacaoPreferencias} \\ \hline
\ttsplit{facilities}    & Gestão de Instalações      & \ttsplit{Espacos}, \ttsplit{Reservas} \\ \hline
\ttsplit{reports}       & Serviço de Relatórios      & \ttsplit{ResumosFinanceiros}, \ttsplit{ResumosDesportivos} \\ \hline
\end{xltabular}
}
\rowcolors{2}{white}{white}
\end{center}

A atribuição do esquema é feita no \texttt{DbContext} de cada módulo com \texttt{modelBuilder.HasDefaultSchema("nome")}. Abaixo, um excerto do \texttt{UsersDbContext}:

\begin{lstlisting}[language={[Sharp]C},basicstyle=\ttfamily\small,breaklines=true,frame=single,caption={Excerto do \texttt{UsersDbContext} com configuração de esquema e herança TPH.},label={lst:usersdbcontext}]
public class UsersDbContext : DbContext
{
    public DbSet<Utilizador> Utilizadores => Set<Utilizador>();
    public DbSet<Atleta> Atletas => Set<Atleta>();
    public DbSet<Treinador> Treinadores => Set<Treinador>();
    public DbSet<EncarregadoEducacao> EncarregadosEducacao
        => Set<EncarregadoEducacao>();

    protected override void OnModelCreating(ModelBuilder mb)
    {
        mb.HasDefaultSchema("users");

        mb.Entity<Utilizador>()
            .HasDiscriminator<string>("UserType")
            .HasValue<Utilizador>("Base")
            .HasValue<Atleta>("Atleta")
            .HasValue<Treinador>("Treinador")
            .HasValue<EncarregadoEducacao>("EncarregadoEducacao");
    }
}
\end{lstlisting}

O \texttt{HasDefaultSchema("users")} mantém as tabelas isoladas dos outros módulos. O \textbf{discriminador TPH} armazena toda a hierarquia de utilizadores numa única tabela física, com uma coluna \texttt{UserType} para distinguir cada tipo. Sacrifica-se alguma normalização em troca de consultas mais simples.

\subsection{Entidades e Modelo de Dados}

Todas as entidades herdam de \texttt{Entity}, que fornece um identificador \texttt{Guid} gerado automaticamente:

\begin{lstlisting}[language={[Sharp]C},basicstyle=\ttfamily\small,breaklines=true,frame=single,caption={Classe base \texttt{Entity} partilhada por todas as entidades do sistema.},label={lst:entity-base}]
public abstract class Entity
{
    public Guid Id { get; protected set; }
    protected Entity() => Id = Guid.NewGuid();
}
\end{lstlisting}

O uso de \texttt{Guid} como chave primária facilita a geração de identificadores no lado do cliente e evita conflitos em inserções concorrentes. Abaixo, a entidade \texttt{Utilizador}:

\begin{lstlisting}[language={[Sharp]C},basicstyle=\ttfamily\small,breaklines=true,frame=single,caption={Classe \texttt{Utilizador} --- entidade base da hierarquia de utilizadores.},label={lst:utilizador}]
public class Utilizador : Entity
{
    public string NumeroSocio { get; set; }
    public string Nome { get; set; }
    public string Email { get; set; }
    public string Telefone { get; set; }
    public string PasswordHash { get; set; }
    public string Role { get; set; }  // Atleta, Treinador, Secretaria, Gerencia
    public string Tipo { get; set; }  // Regular, Jovem, Honorario, Fundador
    public string Estado { get; set; } // Ativo, Suspenso
    public DateTime DataInscricao { get; set; }
    public int FailedLoginAttempts { get; set; }
    public DateTime? LockoutEnd { get; set; }
}
\end{lstlisting}

Os campos \texttt{FailedLoginAttempts} e \texttt{LockoutEnd} suportam a proteção contra força bruta (\textbf{RNF-01}). A classe \texttt{Atleta} estende \texttt{Utilizador} com dados desportivos:

\begin{lstlisting}[language={[Sharp]C},basicstyle=\ttfamily\small,breaklines=true,frame=single,caption={Classe \texttt{Atleta} --- extensão de \texttt{Utilizador} com dados desportivos.},label={lst:atleta}]
public class Atleta : Utilizador
{
    public string Posicao { get; set; }
    public int NumeroCamisola { get; set; }
    public DateTime DataNascimento { get; set; }
    public Guid? EncarregadoEducacaoId { get; set; }
    public EncarregadoEducacao? EncarregadoEducacao { get; set; }
}
\end{lstlisting}

A propriedade de navegação \texttt{EncarregadoEducacao} estabelece a relação com o responsável legal, traduzida numa chave estrangeira dentro da mesma tabela.

A entidade \texttt{Treino} ilustra uma relação de composição: cada treino pertence a uma equipa e tem uma coleção de presenças:

\begin{lstlisting}[language={[Sharp]C},basicstyle=\ttfamily\small,breaklines=true,frame=single,caption={Entidades \texttt{Treino} e \texttt{PresencaTreino} do módulo desportivo.},label={lst:treino}]
public class Treino : Entity
{
    public Guid EquipaId { get; set; }
    public Equipa? Equipa { get; set; }
    public DateTime DataInicio { get; set; }
    public DateTime DataFim { get; set; }
    public Guid EspacoId { get; set; }
    public string Descricao { get; set; }
    public ICollection<PresencaTreino> Presencas { get; set; }
}

public class PresencaTreino : Entity
{
    public Guid TreinoId { get; set; }
    public Guid AtletaId { get; set; }
    public EstadoPresencaTreino Estado { get; set; }
    public string Justificacao { get; set; }
}
\end{lstlisting}

O \texttt{enum EstadoPresencaTreino} tem cinco estados: \texttt{Pendente}, \texttt{Presente}, \texttt{FaltaJustificada}, \texttt{FaltaInjustificada} e \texttt{AusentePorLesao}.

\subsection{Mecanismo de Migrações}

O Entity Framework Core usa \textit{Migrations} para gerir a evolução do esquema da base de dados. Quando uma classe é alterada, é gerada uma nova migração com as instruções correspondentes, versionada no Git.

A aplicação automática das migrações foi restringida a ambientes de desenvolvimento, para evitar que uma instância de produção aplique alterações irreversíveis no arranque.

Abaixo, um excerto da migração inicial do módulo de Utilizadores:

\begin{lstlisting}[language={[Sharp]C},basicstyle=\ttfamily\small,breaklines=true,frame=single,caption={Excerto da migração \texttt{InitialUsers} --- criação da tabela \texttt{Utilizadores}.},label={lst:migration}]
protected override void Up(MigrationBuilder migrationBuilder)
{
    migrationBuilder.EnsureSchema(name: "users");

    migrationBuilder.CreateTable(
        name: "Utilizadores",
        schema: "users",
        columns: table => new
        {
            Id        = table.Column<Guid>(type: "uuid", nullable: false),
            Nome      = table.Column<string>(type: "text", nullable: false),
            Email     = table.Column<string>(type: "text", nullable: false),
            PasswordHash = table.Column<string>(type: "text", nullable: false),
            Role      = table.Column<string>(type: "text", nullable: false),
            UserType  = table.Column<string>(type: "text", nullable: false),
            // ... campos adicionais omitidos por brevidade
        },
        constraints: table =>
        {
            table.PrimaryKey("PK_Utilizadores", x => x.Id);
        });
}
\end{lstlisting}

O programador nunca escreveu SQL. O Entity Framework Core gera as instruções DDL a partir das classes C\#, e \texttt{EnsureSchema} garante que o esquema é criado antes da tabela.

As migrações são aplicadas no arranque via \texttt{Database.Migrate()} no \texttt{Program.cs}:

\begin{lstlisting}[language={[Sharp]C},basicstyle=\ttfamily\small,breaklines=true,frame=single,caption={Aplicação automática das migrações no arranque da API.},label={lst:migrate}]
usersDb.Database.Migrate();
sportsDb.Database.Migrate();
clinicalDb.Database.Migrate();
financeDb.Database.Migrate();
notificationsDb.Database.Migrate();
facilitiesDb.Database.Migrate();
reportsDb.Database.Migrate();
\end{lstlisting}

Esta abordagem garante que a base de dados está sempre na versão mais recente antes de servir pedidos. Cada migração inclui também um método \texttt{Down} para reverter alterações.

\subsection{Povoamento da Base de Dados}

Foi desenvolvido um \texttt{DataSeeder} que povoa automaticamente a base de dados no arranque, caso esteja vazia. Verifica a existência de um registo sentinela para evitar duplicações:

\begin{lstlisting}[language={[Sharp]C},basicstyle=\ttfamily\small,breaklines=true,frame=single,caption={Verificação de idempotência no \texttt{DataSeeder}.},label={lst:seeder-check}]
public static async Task SeedAsync(
    UsersDbContext usersDb,
    SportsDbContext sportsDb,
    ClinicalDbContext clinicalDb,
    FinanceDbContext financeDb,
    FacilitiesDbContext facilitiesDb)
{
    if (usersDb.Utilizadores.Any(u => u.Email == "secretaria@cap.pt"))
        return; // BD ja povoada --- nao duplicar dados
    // ... restante logica de povoamento
}
\end{lstlisting}

O processo segue a ordem de dependências: primeiro os utilizadores, depois as equipas, os treinos e por fim os dados financeiros e clínicos.

Os dados gerados pelo \textit{seeder} foram pensados para refletir o ambiente real de um clube desportivo:

\begin{itemize}
    \item \textbf{Utilizadores}: São criados 2 funcionários (secretaria e gerência), 8 treinadores (um por equipa), 12 encarregados de educação e dezenas de atletas distribuídos por 8 equipas de 3 modalidades (Futebol, Basquetebol e Natação).
    \item \textbf{Treinos e presenças}: Para cada equipa, são gerados 10 treinos passados com intervalos de 3 dias. As presenças de cada atleta são atribuídas de forma variada com recurso a um gerador pseudo-aleatório, garantindo que diferentes atletas apresentam diferentes taxas de assiduidade.
    \item \textbf{Atestados médicos}: São criados atestados com datas de emissão e expiração variadas, simulando cenários de atletas com documentação válida, prestes a expirar e já expirada.
    \item \textbf{Quotas e pagamentos}: São geradas definições de quotas mensais e anuais, com pagamentos em diferentes estados (pago, pendente, em atraso), permitindo testar os alertas do módulo financeiro.
    \item \textbf{Instalações}: São criados 6 espaços físicos (campos, pavilhão, piscina) com reservas distribuídas ao longo da semana.
\end{itemize}

O povoamento programático em C\# traz duas vantagens: os dados são gerados relativamente à data atual (os treinos e atestados estão sempre ``frescos''), e as presenças são geradas com valores aleatórios, evitando dados uniformes e irrealistas.

% ============================================================
% SECÇÃO 5.3 — DESENVOLVIMENTO DA CAMADA DE DADOS
% ============================================================

\section{Desenvolvimento da Camada de Dados}

Com a base de dados implementada, construiu-se a camada de dados, responsável por abstrair o acesso às tabelas. Esta camada usa o \textbf{Entity Framework Core} com o \textbf{padrão Repository} \cite{fowler2002}, sem necessidade de escrever SQL.

\subsection{Interface Genérica de Repositório}

Foi definida uma interface genérica \texttt{IRepository<T>} no projeto partilhado \texttt{CAP.Shared}, com as operações CRUD fundamentais:

\begin{lstlisting}[language={[Sharp]C},basicstyle=\ttfamily\small,breaklines=true,frame=single,caption={Interface genérica \texttt{IRepository<T>} partilhada por todos os módulos.},label={lst:irepository}]
public interface IRepository<T> where T : Entity
{
    Task<T?> GetByIdAsync(Guid id);
    Task<IEnumerable<T>> GetAllAsync();
    Task AddAsync(T entity);
    Task UpdateAsync(T entity);
    Task DeleteAsync(T entity);
    Task SaveChangesAsync();
}
\end{lstlisting}

A restrição \texttt{where T : Entity} garante que apenas entidades com \texttt{Guid} podem ser geridas. Todas as operações são assíncronas para não bloquear \textit{threads}.

\subsection{Implementação Genérica com Entity Framework Core}

Cada módulo tem a sua própria implementação concreta do repositório. Abaixo, o repositório genérico do módulo desportivo:

\begin{lstlisting}[language={[Sharp]C},basicstyle=\ttfamily\small,breaklines=true,frame=single,caption={Implementação genérica \texttt{SportsRepository<T>} do módulo desportivo.},label={lst:sportsrepo}]
public class SportsRepository<T> : IRepository<T> where T : Entity
{
    private readonly SportsDbContext _context;

    public SportsRepository(SportsDbContext context)
        => _context = context;

    public async Task<T?> GetByIdAsync(Guid id)
        => await _context.Set<T>().FindAsync(id);

    public async Task<IEnumerable<T>> GetAllAsync()
        => await _context.Set<T>().ToListAsync();

    public async Task AddAsync(T entity)
        => await _context.Set<T>().AddAsync(entity);

    public Task UpdateAsync(T entity)
    {
        _context.Set<T>().Update(entity);
        return Task.CompletedTask;
    }

    public Task DeleteAsync(T entity)
    {
        _context.Set<T>().Remove(entity);
        return Task.CompletedTask;
    }

    public async Task SaveChangesAsync()
        => await _context.SaveChangesAsync();
}
\end{lstlisting}

O método \texttt{Set<T>()} devolve o \texttt{DbSet} do tipo correspondente, permitindo que um único repositório sirva todas as entidades do módulo. O Entity Framework traduz os métodos como \texttt{FindAsync} ou \texttt{AddAsync} nas instruções SQL correspondentes.

\subsection{Repositórios Especializados}

Alguns módulos precisam de métodos de consulta adicionais. O repositório de Utilizadores tem três métodos especializados:

\begin{lstlisting}[language={[Sharp]C},basicstyle=\ttfamily\small,breaklines=true,frame=single,caption={Métodos especializados do \texttt{UserRepository}.},label={lst:userrepo}]
public async Task<Utilizador?> GetByEmailAsync(string email)
    => await _context.Utilizadores
        .FirstOrDefaultAsync(u => u.Email == email);

public async Task<EncarregadoEducacao?> GetEncarregadoWithDependentesAsync(
    Guid id)
    => await _context.Utilizadores
        .OfType<EncarregadoEducacao>()
        .Include(e => e.AtletasDependentes)
        .FirstOrDefaultAsync(e => e.Id == id);

public async Task<Atleta?> GetAtletaByIdAsync(Guid id)
    => await _context.Utilizadores
        .OfType<Atleta>()
        .FirstOrDefaultAsync(a => a.Id == id);
\end{lstlisting}

O \texttt{OfType<T>()} filtra pelo tipo concreto na hierarquia TPH. O \texttt{Include()} carrega propriedades de navegação numa única consulta, evitando o problema de \textit{N+1 queries}.

\subsection{Injeção de Dependências}

Os repositórios são registados com ciclo de vida \texttt{Scoped}. A configuração foi extraída para \texttt{ModuleExtensions.cs} usando \textit{Extension Methods}, para evitar um \texttt{Program.cs} demasiado grande:

\begin{lstlisting}[language={[Sharp]C},basicstyle=\ttfamily\small,breaklines=true,frame=single,caption={Registo delegado aos métodos de extensão (\texttt{Program.cs}).},label={lst:di}]
// Dependency Injection via Extension Methods
builder.Services.AddUsersModule()
                .AddSportsModule()
                .AddClinicalModule()
                .AddFinanceModule()
                .AddNotificationsModule()
                .AddFacilitiesModule()
                .AddReportsModule();
\end{lstlisting}

O registo de tipo aberto (\textit{open generic}) evita registar manualmente cada combinação de repositório — uma única linha resolve todos os pedidos.

\subsection{Resumo da Arquitetura da Camada de Dados}

A tabela seguinte resume a relação entre módulos, contextos e repositórios:

\begin{center}
{\footnotesize
\rowcolors{2}{rowblue}{white}
\begin{xltabular}{\textwidth}{|>{\raggedright\arraybackslash}p{2cm}|>{\raggedright\arraybackslash}p{3.7cm}|>{\raggedright\arraybackslash}p{4.2cm}|>{\raggedright\arraybackslash}X|}
\hline
\rowcolor{headerblue}\textcolor{white}{\textbf{Módulo}} & \textcolor{white}{\textbf{DbContext}} & \textcolor{white}{\textbf{Repositório}} & \textcolor{white}{\textbf{Entidades}} \\ \hline
\endfirsthead

\multicolumn{4}{c}{{\tablename\ \thetable{} -- Continuação}} \\
\hline
\rowcolor{headerblue}\textcolor{white}{\textbf{Módulo}} & \textcolor{white}{\textbf{DbContext}} & \textcolor{white}{\textbf{Repositório}} & \textcolor{white}{\textbf{Entidades}} \\ \hline
\endhead

\hline \multicolumn{4}{|r|}{\textit{Continua na próxima página\ldots}} \\ \hline \endfoot

\hline
\caption{Mapeamento entre módulos, contextos de dados e repositórios.}
\label{tab:camada-dados} \\
\endlastfoot

Utilizadores  & \ttsplit{UsersDbContext}         & \ttsplit{UserRepository}                & Utilizador, Atleta, Treinador, Encarregado \\ \hline
Desportivo    & \ttsplit{SportsDbContext}        & \ttsplit{SportsRepository<T>}           & Modalidade, Equipa, Treino, Convocatoria, Presença \\ \hline
Clínico       & \ttsplit{ClinicalDbContext}      & \ttsplit{ClinicalRepository<T>}         & AtestadoMedico, Lesao \\ \hline
Financeiro    & \ttsplit{FinanceDbContext}       & \ttsplit{FinanceRepository<T>}          & QuotaDefinicao, Quota, Pagamento, Recibo \\ \hline
Notificações  & \ttsplit{NotificationsDbContext} & \ttsplit{NotificationsRepository<T>}    & Notificacao, NotificacaoPreferencia \\ \hline
Instalações   & \ttsplit{FacilitiesDbContext}    & \ttsplit{FacilitiesRepository<T>}       & Espaco, Reserva \\ \hline
Relatórios    & \ttsplit{ReportsDbContext}       & \ttsplit{ReportsRepository<T>}          & ResumoFinanceiro, ResumoDesportivo \\ \hline
\end{xltabular}
}
\rowcolors{2}{white}{white}
\end{center}

Esta organização garante que cada módulo acede apenas ao seu esquema. A comunicação entre módulos é feita na camada de negócio, nunca na camada de dados.

% ============================================================
% SECÇÃO 5.4 — DESENVOLVIMENTO DA CAMADA DE NEGÓCIO
% ============================================================

\section{Desenvolvimento da Camada de Negócio}\label{sec:rel_backend}
\verIA{subsec:ia_impl_backend}

A camada de negócio corresponde aos \textbf{Controllers} da API REST. Cada controlador recebe pedidos HTTP, aplica regras de negócio e usa os repositórios para interagir com a base de dados.

\subsection{Autenticação com JSON Web Tokens}

A autenticação usa \textbf{JWT} conforme o \textbf{RNF-03}. No \textit{login}, o servidor gera um \textit{token} com o identificador, nome, email e papel do utilizador. O \textit{frontend} envia este \textit{token} em cada pedido no cabeçalho \texttt{Authorization}.

\begin{lstlisting}[language={[Sharp]C},basicstyle=\ttfamily\small,breaklines=true,frame=single,caption={Geração de \textit{token} JWT com as \textit{claims} do utilizador.},label={lst:jwt}]
public string GenerateToken(Utilizador user)
{
    var key = Encoding.ASCII.GetBytes(
        _configuration["Jwt:Key"]);

    var tokenDescriptor = new SecurityTokenDescriptor
    {
        Subject = new ClaimsIdentity(new[]
        {
            new Claim(ClaimTypes.NameIdentifier, user.Id.ToString()),
            new Claim(ClaimTypes.Name, user.Nome),
            new Claim(ClaimTypes.Email, user.Email),
            new Claim(ClaimTypes.Role, user.Role)
        }),
        Expires = DateTime.UtcNow.AddDays(7),
        SigningCredentials = new SigningCredentials(
            new SymmetricSecurityKey(key),
            SecurityAlgorithms.HmacSha256Signature)
    };

    var token = tokenHandler.CreateToken(tokenDescriptor);
    return tokenHandler.WriteToken(token);
}
\end{lstlisting}

O \textit{token} é válido por 7 dias e assinado com HMAC-SHA256. O \texttt{ClaimTypes.Role} permite o controlo de acessos por papel.

\subsection{Hashing e Segurança de Passwords}

As \textit{passwords} nunca são armazenadas em texto simples. A biblioteca \texttt{BCrypt.Net-Next} gera uma \textit{hash} no registo, e no \textit{login} invoca \texttt{BCrypt.Verify()} para comparar com o valor armazenado. O \texttt{DataSeeder} também usa \textit{hashing} nos utilizadores predefinidos.

\subsection{Proteção contra Ataques de Força Bruta}

O \texttt{AuthController} implementa a proteção contra força bruta (\textbf{RNF-01}). Após 5 tentativas falhadas, a conta é bloqueada 15 minutos:

\begin{lstlisting}[language={[Sharp]C},basicstyle=\ttfamily\small,breaklines=true,frame=single,caption={Mecanismo de bloqueio por tentativas de \textit{login} falhadas.},label={lst:lockout}]
[HttpPost("login")]
public async Task<IActionResult> Login(LoginRequest request)
{
    var user = await _userRepository.GetByEmailAsync(request.Email);
    if (user == null) return Unauthorized("Credenciais invalidas");

    // Verificar se a conta esta bloqueada
    if (user.LockoutEnd.HasValue
        && user.LockoutEnd.Value > DateTime.UtcNow)
        return Unauthorized("Conta bloqueada.");

    if (!BCrypt.Net.BCrypt.Verify(request.Password, user.PasswordHash))
    {
        user.FailedLoginAttempts++;
        if (user.FailedLoginAttempts >= 5)
        {
            user.LockoutEnd = DateTime.UtcNow.AddMinutes(15);
            user.FailedLoginAttempts = 0;
        }
        await _userRepository.UpdateAsync(user);
        await _userRepository.SaveChangesAsync();
        return Unauthorized("Credenciais invalidas");
    }

    // Login bem-sucedido: resetar contadores
    user.FailedLoginAttempts = 0;
    user.LockoutEnd = null;
    // ... gerar token JWT
}
\end{lstlisting}

No \textit{login} com sucesso, os contadores são reiniciados. Este mecanismo impede ataques por tentativa exaustiva.

\subsection{Controlo de Acessos por Papel (RBAC)}

O controlo de acessos é implementado com o atributo \texttt{[Authorize]} do ASP.NET Core. Cada \textit{endpoint} especifica os papéis permitidos e o \textit{middleware} valida automaticamente o token JWT.

\begin{lstlisting}[language={[Sharp]C},basicstyle=\ttfamily\small,breaklines=true,frame=single,caption={Exemplos de controlo de acessos RBAC nos \textit{Controllers}.},label={lst:rbac}]
// Apenas Treinadores e Gerencia podem criar treinos
[HttpPost]
[Authorize(Roles = "Treinador,Gerencia")]
public async Task<IActionResult> Create(CreateTreinoRequest request)
{ ... }

// Apenas Treinadores podem registar presencas
[HttpPost("{id}/attendance")]
[Authorize(Roles = "Treinador")]
public async Task<IActionResult> RecordAttendance(
    Guid id, List<UpdatePresencaRequest> requests)
{ ... }

// Atletas, Treinadores e Gerencia podem consultar treinos
[HttpGet]
[Authorize(Roles = "Treinador,Gerencia,Atleta")]
public async Task<IActionResult> GetAll()
{ ... }
\end{lstlisting}

As permissões da Matriz RBAC são aplicadas diretamente ao nível de cada rota. Um atleta que tente criar um treino recebe automaticamente HTTP 403.

\subsection{Estrutura dos Controladores por Módulo}

A tabela seguinte lista os controladores implementados por módulo:

\begin{center}
{\footnotesize
\rowcolors{2}{rowblue}{white}
\begin{xltabular}{\textwidth}{|>{\raggedright\arraybackslash}p{2cm}|>{\raggedright\arraybackslash}p{4cm}|>{\raggedright\arraybackslash}p{3.7cm}|>{\raggedright\arraybackslash}X|}
\hline
\rowcolor{headerblue}\textcolor{white}{\textbf{Módulo}} & \textcolor{white}{\textbf{Controlador}} & \textcolor{white}{\textbf{Rota Base}} & \textcolor{white}{\textbf{Funcionalidades}} \\ \hline
\endfirsthead

\multicolumn{4}{c}{{\tablename\ \thetable{} -- Continuação}} \\
\hline
\rowcolor{headerblue}\textcolor{white}{\textbf{Módulo}} & \textcolor{white}{\textbf{Controlador}} & \textcolor{white}{\textbf{Rota Base}} & \textcolor{white}{\textbf{Funcionalidades}} \\ \hline
\endhead

\hline \multicolumn{4}{|r|}{\textit{Continua na próxima página\ldots}} \\ \hline \endfoot

\hline
\caption{Inventário dos controladores da API REST, organizados por módulo.}
\label{tab:controladores} \\
\endlastfoot

Utilizadores  & \ttsplit{AuthController}          & \ttsplit{/api/users/auth}         & Login, Registo \\ \hline
Utilizadores  & \ttsplit{AthletesController}      & \ttsplit{/api/users/athletes}     & Listar e consultar atletas \\ \hline
Utilizadores  & \ttsplit{ProfileController}       & \ttsplit{/api/users/profile}      & Editar perfil próprio \\ \hline
Utilizadores  & \ttsplit{ParentalController}      & \ttsplit{/api/users/parental}     & Relação EE $\leftrightarrow$ Atleta \\ \hline
Utilizadores  & \ttsplit{SociosController}        & \ttsplit{/api/users/socios}       & Gestão de sócios \\ \hline
Desportivo    & \ttsplit{TeamsController}         & \ttsplit{/api/sports/teams}       & Equipas e plantéis \\ \hline
Desportivo    & \ttsplit{TreinosController}       & \ttsplit{/api/sports/trainings}   & Treinos e presenças \\ \hline
Desportivo    & \ttsplit{ConvocatoriasController} & \ttsplit{/api/sports/calls}       & Convocatórias e convites \\ \hline
Desportivo    & \ttsplit{OcorrenciasController}   & \ttsplit{/api/sports/events}      & Resultados de jogos \\ \hline
Clínico       & \ttsplit{AtestadosController}     & \ttsplit{/api/clinical/certs}     & Submissão e validação de atestados \\ \hline
Clínico       & \ttsplit{LesoesController}        & \ttsplit{/api/clinical/injuries}  & Reportar e gerir lesões \\ \hline
Financeiro    & \ttsplit{QuotasController}        & \ttsplit{/api/finance/quotas}     & Definição e consulta de quotas \\ \hline
Financeiro    & \ttsplit{PagamentosController}    & \ttsplit{/api/finance/payments}   & Registo de pagamentos \\ \hline
Financeiro    & \ttsplit{WebhooksController}      & \ttsplit{/api/finance/webhooks}   & Confirmação assíncrona \\ \hline
Notificações  & \ttsplit{NotificationsController} & \ttsplit{/api/notifications}      & Envio e consulta de alertas \\ \hline
Instalações   & \ttsplit{FacilitiesController}    & \ttsplit{/api/facilities}         & Espaços e reservas \\ \hline
Instalações   & \ttsplit{LojaController}          & \ttsplit{/api/facilities/loja}    & Catálogo e encomendas de fardamento \\ \hline
Relatórios    & \ttsplit{ReportsController}       & \ttsplit{/api/reports}            & Exportação SAF-T e resumos \\ \hline
\end{xltabular}
}
\rowcolors{2}{white}{white}
\end{center}

\subsection{Exemplo de Regra de Negócio: Registo de Presenças}

Para ilustrar a lógica de negócio, apresenta-se o fluxo de registo de presenças. O controlador recebe uma lista de pares (\texttt{AtletaId}, \texttt{Estado}) e atualiza os registos:

\begin{lstlisting}[language={[Sharp]C},basicstyle=\ttfamily\small,breaklines=true,frame=single,caption={Registo de presenças com validação de treino existente.},label={lst:attendance}]
[HttpPost("{id}/attendance")]
[Authorize(Roles = "Treinador")]
public async Task<IActionResult> RecordAttendance(
    Guid id, List<UpdatePresencaRequest> requests)
{
    var treino = await _dbContext.Treinos
        .Include(t => t.Presencas)
        .FirstOrDefaultAsync(t => t.Id == id);
    if (treino == null) return NotFound("Treino nao encontrado");

    foreach (var req in requests)
    {
        var presenca = treino.Presencas
            .FirstOrDefault(p => p.AtletaId == req.AtletaId);
        if (presenca != null)
        {
            presenca.Estado = req.Estado;
            presenca.Justificacao = req.Justificacao ?? "";
        }
        else
        {
            treino.Presencas.Add(new PresencaTreino
            {
                TreinoId = id,
                AtletaId = req.AtletaId,
                Estado = req.Estado,
                Justificacao = req.Justificacao ?? ""
            });
        }
    }
    await _dbContext.SaveChangesAsync();
    return Ok("Assiduidade registada com sucesso");
}
\end{lstlisting}

O \texttt{Include()} carrega as presenças numa única consulta. A lógica segue um padrão \textit{upsert}: atualiza se já existe, cria se não existe. O \texttt{SaveChangesAsync()} persiste tudo numa única transação.

\subsection{Comunicação Inter-Módulo com MediatR}

Para a comunicação entre módulos foi configurado o \textbf{MediatR}, uma biblioteca que implementa o padrão \textit{Mediator} em memória:

\begin{lstlisting}[language={[Sharp]C},basicstyle=\ttfamily\small,breaklines=true,frame=single,caption={Configuração do MediatR para comunicação inter-módulo.},label={lst:mediatr}]
builder.Services.AddMediatR(cfg =>
    cfg.RegisterServicesFromAssemblies(
        AppDomain.CurrentDomain.GetAssemblies()));
\end{lstlisting}

Esta linha regista todos os \textit{handlers} de eventos dos módulos. Na prática, a comunicação entre módulos ainda é feita maioritariamente por chamadas diretas entre repositórios. A adoção de eventos é um passo previsto para o futuro.

% ============================================================
% SECÇÃO 5.5 — DESENVOLVIMENTO DA CAMADA DE APRESENTAÇÃO
% ============================================================

\section{Desenvolvimento da Camada de Apresentação}\label{sec:rel_frontend}
\verIA{subsec:ia_impl_frontend}

A camada de apresentação foi construída com \textbf{Next.js~16} (React~19, TypeScript), conforme o \textbf{RNF-05} \cite{nextjs2024}. As subsecções seguintes descrevem os aspetos mais relevantes desta camada.

\subsection{Componentes Reutilizáveis}

Foi construída uma biblioteca de componentes reutilizáveis baseada no \textbf{shadcn/ui} \cite{shadcnui2024}. A tabela seguinte resume os componentes implementados:

\begin{center}
{\small
\rowcolors{2}{rowblue}{white}
\begin{xltabular}{\textwidth}{|>{\raggedright\arraybackslash}p{3.5cm}|>{\raggedright\arraybackslash}X|}
\hline
\rowcolor{headerblue}\textcolor{white}{\textbf{Componente}} & \textcolor{white}{\textbf{Descrição}} \\ \hline
\endfirsthead

\multicolumn{2}{c}{{\tablename\ \thetable{} -- Continuação}} \\
\hline
\rowcolor{headerblue}\textcolor{white}{\textbf{Componente}} & \textcolor{white}{\textbf{Descrição}} \\ \hline
\endhead

\hline \multicolumn{2}{|r|}{\textit{Continua na próxima página\ldots}} \\ \hline \endfoot

\hline
\caption{Componentes reutilizáveis da interface, baseados na biblioteca shadcn/ui.}
\label{tab:componentes-ui} \\
\endlastfoot

\texttt{Button}          & Botão com variantes (primário, secundário, destrutivo) e estados de carregamento. \\ \hline
\texttt{Card}            & Contentor com cabeçalho, conteúdo e sombra, utilizado nos \textit{dashboards} e formulários. \\ \hline
\texttt{Input / Label}   & Campos de entrada com validação e acessibilidade. \\ \hline
\texttt{Table}           & Tabela estilizada para listagens de atletas, treinos, quotas, etc. \\ \hline
\texttt{Dialog / Sheet}  & Modais e painéis laterais para ações sem sair da página (criar treino, editar atleta). \\ \hline
\texttt{Select / Tabs}   & Seleção de opções e separadores para organizar conteúdo em abas. \\ \hline
\texttt{Calendar}        & Componente de calendário visual para visualização de treinos e convocatórias. \\ \hline
\texttt{Badge / Alert}   & Indicadores de estado (``Pago'', ``Em atraso'') e mensagens contextuais. \\ \hline
\texttt{Sidebar}         & Barra de navegação lateral com ícones, responsiva e colapsável. \\ \hline
\texttt{Avatar / Tooltip} & Fotografia do utilizador e dicas de contexto ao passar o rato. \\ \hline
\end{xltabular}
}
\rowcolors{2}{white}{white}
\end{center}

\subsection{Gestão de Sessão e Comunicação com a API}

A sessão é guardada no \texttt{localStorage} do navegador com o \textit{token} JWT e os dados do utilizador. Todos os pedidos à API incluem o \textit{token} no cabeçalho \texttt{Authorization}.

Para centralizar a comunicação com a API e evitar repetição de código, foi criada a função utilitária \texttt{fetchApi}:

\begin{lstlisting}[language=Java,basicstyle=\ttfamily\small,breaklines=true,frame=single,caption={Função utilitária \texttt{fetchApi} para comunicação com a API REST.},label={lst:fetchapi}]
export async function fetchApi<T>(
    endpoint: string, options: RequestInit = {}): Promise<T>
{
    const token = localStorage.getItem('cap_token');

    const headers = new Headers(options.headers);
    if (token) {
        headers.set('Authorization', `Bearer ${token}`);
    }
    if (!headers.has('Content-Type')
        && !(options.body instanceof FormData)) {
        headers.set('Content-Type', 'application/json');
    }

    const response = await fetch(
        `${API_BASE_URL}${endpoint}`, { ...options, headers });

    if (!response.ok) {
        if (response.status === 401) {
            localStorage.removeItem('cap_token');
            window.location.href = '/'; // Redirecionar para login
        }
        throw new Error(`Erro HTTP: ${response.status}`);
    }
    return response.json();
}
\end{lstlisting}

A função injeta o \textit{token} JWT em todos os pedidos, trata o código HTTP 401 redirecionando para o \textit{login}, e suporta \texttt{FormData} para \textit{upload} de ficheiros como atestados em PDF.

\subsection{Navegação e Routing por Papel}

A navegação usa o \textit{App Router} do Next.js. A estrutura de diretórios reflete a hierarquia de URLs:

\begin{lstlisting}[basicstyle=\ttfamily\small,breaklines=true,frame=single,caption={Estrutura de diretórios do \textit{frontend}, organizada por papel de utilizador.},label={lst:routing}]
app/
  page.tsx                  # Login (/)
  dashboard/
    atleta/
      page.tsx              # /dashboard/atleta
      calendario/page.tsx   # /dashboard/atleta/calendario
      convocatorias/page.tsx
      loja/page.tsx
    treinador/
      page.tsx              # /dashboard/treinador
      plantel/page.tsx
      presencas/page.tsx
      convocatorias/page.tsx
      calendario/page.tsx
      atestados/page.tsx
      relatorios/page.tsx
    encarregado/
      page.tsx              # /dashboard/encarregado
      atletas/page.tsx
      quotas/page.tsx
      pagamentos/page.tsx
      atestados/page.tsx
    secretaria/
      page.tsx              # /dashboard/secretaria
      atletas/page.tsx
      socios/page.tsx
      quotas/page.tsx
      instalacoes/page.tsx
      relatorios/page.tsx
    definicoes/page.tsx     # Definicoes (todos os papeis)
    notificacoes/page.tsx   # Notificacoes (todos os papeis)
\end{lstlisting}

Após o \textit{login}, o utilizador é redirecionado para o \textit{dashboard} do seu papel. Cada \textit{dashboard} tem a sua barra de navegação com apenas as páginas relevantes para esse papel.

\subsection{Página de Login e Fluxo de Autenticação}

A página de \textit{login} é o ponto de entrada da aplicação. O formulário envia as credenciais para \texttt{POST /api/users/auth/login}. Em caso de sucesso, o \textit{token} JWT \cite{rfc7519} é guardado no \texttt{localStorage} e o utilizador é redirecionado para o seu \textit{dashboard}:

\begin{lstlisting}[language=Java,basicstyle=\ttfamily\small,breaklines=true,frame=single,caption={Fluxo de \textit{login} no \textit{frontend} React.},label={lst:login-flow}]
const handleLogin = async (e: React.FormEvent) => {
    const response = await fetch(
        '/api/users/auth/login',
        { method: 'POST',
          body: JSON.stringify({ email, password }) });

    if (response.ok) {
        const data = await response.json();
        localStorage.setItem("cap_token", data.token);
        localStorage.setItem("cap_user",
            JSON.stringify({ email, role: data.role, nome: data.nome }));
        router.push(`/dashboard/${data.role.toLowerCase()}`);
    } else {
        setError("Email ou password incorretos");
    }
};
\end{lstlisting}

\subsection{Estrutura de um Dashboard}

Cada \textit{dashboard} usa \texttt{useEffect} para carregar dados da API ao iniciar, e \texttt{useState} para guardar esses dados no estado local. O \textit{dashboard} do Treinador, por exemplo, mostra cinco cartões resumo (atletas, assiduidade, próximo treino, convocatórias, presenças), cada um com a sua própria chamada \texttt{fetchApi}.

O sistema de \textit{dashboards} foi desenhado para que cada papel de utilizador veja apenas a informação que lhe diz respeito:

\begin{itemize}
    \item \textbf{Atleta}: Consulta os seus treinos, presenças, convocatórias e loja de equipamentos.
    \item \textbf{Treinador}: Gere o plantel, regista presenças, cria convocatórias e consulta atestados dos atletas.
    \item \textbf{Encarregado de Educação}: Acompanha os atletas dependentes, consulta quotas e submete atestados.
    \item \textbf{Secretaria}: Gere sócios, atletas, instalações, quotas e pagamentos, e exporta relatórios financeiros.
\end{itemize}

\subsection{Exemplo de Página: Gestão de Presenças}

A página de presenças do treinador ilustra a integração entre \textit{frontend} e \textit{backend}. O fluxo é: (1)~a página carrega os treinos via \texttt{GET /api/sports/trainings}; (2)~o treinador seleciona um treino e vê os atletas com os estados atuais; (3)~escolhe um dos cinco estados por atleta; (4)~ao guardar, os dados são enviados via \texttt{POST /api/sports/trainings/\{id\}/attendance}; (5)~é mostrada uma mensagem de confirmação. Este padrão (carregar $\rightarrow$ formulário $\rightarrow$ submeter $\rightarrow$ confirmar) repete-se em todas as páginas de edição.

% ============================================================
% SECÇÃO 5.6 — AVALIAÇÃO DA APLICAÇÃO DESENVOLVIDA
% ============================================================

\section{Avaliação da Aplicação Desenvolvida}

No final do desenvolvimento, foi feita uma avaliação para verificar se a aplicação cumpre os requisitos do Capítulo~\ref{chap:analise_requisitos} e está conforme a especificação do Capítulo~\ref{chap:especificacao_modelacao}.

\subsection{Verificação dos Requisitos Funcionais}

Cada requisito foi verificado e classificado como: {\large\checkmark} (cumprido), {\large$\sim$} (parcialmente cumprido) ou {\large$\times$} (não cumprido).

{\small
\rowcolors{2}{rowblue}{white}
\begin{xltabular}{\textwidth}{|p{1.5cm}|X|p{1.5cm}|p{3.2cm}|}
\hline
\rowcolor{headerblue}\textcolor{white}{\textbf{Requisito}} & \textcolor{white}{\textbf{Nome}} & \textcolor{white}{\textbf{Estado}} & \textcolor{white}{\textbf{Observações}} \\ \hline
\endfirsthead

\multicolumn{4}{c}{{\tablename\ \thetable{} -- Continuação}} \\
\hline
\rowcolor{headerblue}\textcolor{white}{\textbf{Requisito}} & \textcolor{white}{\textbf{Nome}} & \textcolor{white}{\textbf{Estado}} & \textcolor{white}{\textbf{Observações}} \\ \hline
\endhead

\hline \multicolumn{4}{|r|}{\textit{Continua na próxima página\ldots}} \\ \hline \endfoot

\hline
\caption{Avaliação do cumprimento dos requisitos funcionais.}
\label{tab:avaliacao-rf} \\
\endlastfoot

RF-01  & Gestão de Utilizadores e Perfis             & {\large\checkmark} & CRUD com JWT \\ \hline
RF-02  & Gestão de Perfis Parentais                  & {\large\checkmark} & Vinculação EE--atleta \\ \hline
RF-03  & Sistema de Comunicação Interna (Chat)        & {\large$\times$}   & Adiado (sem chat ainda) \\ \hline
RF-04  & Motor de Notificações Automáticas            & {\large$\sim$}     & Só \textit{in-app} (sem email/SMS) \\ \hline
RF-05  & Gestão de Modalidades e Épocas               & {\large\checkmark} & 3 modalidades ativas \\ \hline
RF-06  & Planeamento de Horários e Treinos            & {\large\checkmark} & Calendário em tempo real \\ \hline
RF-07  & Gestão de Presenças e Faltas                 & {\large\checkmark} & 5 estados de presença \\ \hline
RF-08  & Gestão de Convocatórias e Disponibilidade    & {\large\checkmark} & Com confirmação Sim/Não \\ \hline
RF-09  & Planos de Treino e Recursos Táticos          & {\large$\sim$}     & Apenas \textit{upload}, sem \textit{streaming} \\ \hline
RF-10  & Registo de Ocorrências e Disciplina          & {\large\checkmark} & Golos, cartões, suspensões \\ \hline
RF-11  & Gestão do Exame Médico Desportivo            & {\large\checkmark} & \textit{Upload} + validação \\ \hline
RF-12  & Gestão de Lesões e Recuperação               & {\large\checkmark} & Reportar e bloquear convocatória \\ \hline
RF-13  & Acompanhamento Físico e Evolução de Métricas & {\large\checkmark} & Peso, altura, IMC \\ \hline
RF-14  & Gestão Documental e Federativa               & {\large$\sim$}     & \textit{Upload} sem encriptação dedicada \\ \hline
RF-15  & Painéis de Desempenho (Dashboards)           & {\large\checkmark} & Por papel \\ \hline
RF-16  & Avaliação Qualitativa e Objetivos            & {\large\checkmark} & Notas e objetivos por atleta \\ \hline
RF-17  & Gestão de Quotas e Pagamentos                & {\large$\sim$}     & Registo manual; \textit{gateway} simulado \\ \hline
RF-18  & Contabilidade Global do Clube                & {\large$\sim$}     & Listagem básica de movimentos \\ \hline
RF-19  & Conformidade Fiscal (Exportação SAF-T)       & {\large\checkmark} & Formato XML \\ \hline
RF-20  & Gestão de Espaços Físicos                    & {\large\checkmark} & 6 espaços parametrizados \\ \hline
RF-21  & Manutenção de Infraestruturas                & {\large$\times$}   & Fora do âmbito desta iteração \\ \hline
RF-22  & Gestão de Stock e Loja                       & {\large$\sim$}     & Catálogo, sem checkout real \\ \hline
RF-23  & Gestão de Eventos Extra-desportivos          & {\large$\times$}   & Fora do âmbito desta iteração \\ \hline
\end{xltabular}
}
\rowcolors{2}{white}{white}

Dos 23 requisitos funcionais, 14 foram cumpridos, 6 parcialmente e 3 não cumpridos. Os parciais correspondem a funcionalidades em que a estrutura está implementada mas alguma integração externa ficou por ligar (notificações em RF-04, pagamentos em RF-17, encriptação de documentos em RF-14, vídeo em RF-09, financeiro em RF-18 e loja em RF-22). Os três não cumpridos (chat em RF-03, manutenção em RF-21 e eventos em RF-23) foram excluídos por serem de baixa prioridade.

\subsection{Verificação dos Requisitos Não Funcionais}

Os requisitos não funcionais foram avaliados em ambiente de desenvolvimento. Os relacionados com escalabilidade, \textit{uptime} e \textit{backups} só podem ser verificados em produção.

{\small
\rowcolors{2}{rowblue}{white}
\begin{xltabular}{\textwidth}{|p{1.5cm}|X|p{1.5cm}|p{3.2cm}|}
\hline
\rowcolor{headerblue}\textcolor{white}{\textbf{Requisito}} & \textcolor{white}{\textbf{Nome}} & \textcolor{white}{\textbf{Estado}} & \textcolor{white}{\textbf{Observações}} \\ \hline
\endfirsthead

\multicolumn{4}{c}{{\tablename\ \thetable{} -- Continuação}} \\
\hline
\rowcolor{headerblue}\textcolor{white}{\textbf{Requisito}} & \textcolor{white}{\textbf{Nome}} & \textcolor{white}{\textbf{Estado}} & \textcolor{white}{\textbf{Observações}} \\ \hline
\endhead

\hline \multicolumn{4}{|r|}{\textit{Continua na próxima página\ldots}} \\ \hline \endfoot

\hline
\caption{Avaliação do cumprimento dos requisitos não funcionais.}
\label{tab:avaliacao-rnf} \\
\endlastfoot

RNF-01 & Proteção de Dados Sensíveis (Encriptação)     & {\large\checkmark} & HTTPS local; \textit{password hashing} (bcrypt) \\ \hline
RNF-02 & Conformidade Legal e RGPD                     & {\large$\sim$}     & Estrutura presente, sem auditoria externa \\ \hline
RNF-03 & Controlo de Acessos (RBAC)                    & {\large\checkmark} & \texttt{[Authorize(Roles)]} em todos os \textit{endpoints} \\ \hline
RNF-04 & Design Responsivo (Mobile-First)              & {\large\checkmark} & Next.js + Tailwind CSS \\ \hline
RNF-05 & Interface Intuitiva e Acessibilidade          & {\large\checkmark} & shadcn/ui, \textit{feedback} imediato \\ \hline
RNF-06 & Tempos de Resposta Rápidos (Performance)      & {\large$\sim$}     & Validado só em ambiente local \\ \hline
RNF-07 & Suporte a Picos de Acesso (Escalabilidade)    & {\large$\times$}   & Não testado fora do âmbito local \\ \hline
RNF-08 & Alta Disponibilidade (Uptime)                 & {\large$\times$}   & Sem ambiente de produção \\ \hline
RNF-09 & Cópias de Segurança (Backups)                 & {\large$\times$}   & Política definida, sem execução automática \\ \hline
RNF-10 & Integração com Gateways de Pagamento          & {\large$\times$}   & Simulado; sem Ifthenpay/Easypay ligado \\ \hline
RNF-11 & Exportação de Dados e Normalização             & {\large\checkmark} & SAF-T XML, PDF, Excel \\ \hline
RNF-12 & Arquitetura Modular e Escalável               & {\large\checkmark} & Monólito Modular com 7 esquemas \\ \hline
RNF-13 & Registos de Auditoria (\textit{Logs})         & {\large$\sim$}     & Apenas \textit{logs} de login implementados \\ \hline
\end{xltabular}
}
\rowcolors{2}{white}{white}

Os RNF-07, RNF-08 e RNF-09 só são avaliáveis em produção, por isso foram marcados como não cumpridos.

O \textbf{RNF-01} está cumprido: HTTPS configurado no Docker Compose e \textit{passwords} encriptadas com \texttt{BCrypt.Net-Next}. O RBAC (\textbf{RNF-03}) está implementado com \texttt{[Authorize(Roles)]} em todos os \textit{endpoints}.

Quanto ao \textbf{RNF-06}, em ambiente local com $\approx$\,150 atletas e 80 treinos, as chamadas à API ficaram abaixo de 500\,ms e os \textit{dashboards} carregam em menos de 3 segundos. Estes números não substituem testes de carga reais.

\subsection{Conformidade com a Especificação}

Também foi verificado se a implementação está conforme a especificação do Capítulo~\ref{chap:especificacao_modelacao}:

\begin{itemize}
    \item \textbf{Modelo de domínio}: As entidades e relações foram implementadas nas classes C\#. A geração de PDF da \texttt{Fatura} não foi implementada nesta versão.

    \item \textbf{Arquitetura em camadas}: A separação em camada de dados, camada de negócio e camada de apresentação foi mantida conforme especificado.

    \item \textbf{Diagrama de componentes}: Os 7 módulos correspondem aos componentes especificados. A principal diferença é que a comunicação inter-módulo via MediatR ainda não é usada de forma extensiva.

    \item \textbf{Casos de uso}: A maioria dos casos de uso foi implementada, com validações adicionais não modeladas inicialmente (duplicação de emails, validação de datas). Os casos de uso ligados aos RF excluídos ficam por implementar.

    \item \textbf{Acessibilidade móvel}: O \textit{frontend} foi construído de forma responsiva, acessível em navegadores móveis.

    \item \textbf{Matriz RBAC}: Implementada com \texttt{[Authorize(Roles = "...")]} em conformidade com a Tabela~\ref{tab:rbac}.
\end{itemize}

Desta análise resultaram duas melhorias identificadas para implementar no futuro:

\begin{itemize}
    \item \textbf{Notificações externas}: ligar as notificações a um \textit{gateway} de email/SMS.
    \item \textbf{Paginação}: introduzir paginação nas listagens de atletas, treinos e pagamentos.
\end{itemize}

\subsection{Garantia de Qualidade e Testes Unitários}

A aplicação inclui testes unitários na suíte \texttt{CAP.Modules.Users.Api.Tests}, desenvolvida com xUnit, Moq e o provedor \texttt{EF Core In-Memory Database}. Os testes cobrem o \texttt{AuthController}, verificando o \textit{login} com credenciais corretas e as respostas 401 com credenciais erradas.

Em resumo, a versão atual cumpre 14 dos 23 RF, tem 6 parciais e 3 fora de âmbito. O sistema está operacional para demonstração e a segurança dos dados está garantida.
